\begin{table}[t]
\centering
\caption{Data structures used by MXA*-ID and why.}
\label{tab:structures}
\footnotesize
\begin{tabular}{p{1.45cm}p{1.6cm}p{3.4cm}}
\toprule
Component & Structure & Role and complexity \\
\midrule
Graph & adjacency list & $O(1)$ neighbour iteration; $O(V+E)$ memory. \\
OPEN & indexed binary heap & one entry/node; $O(\log V)$ push, update-key, pop; deterministic keys. \\
Search state & dense arrays & $O(1)$ $g$, parent, closed and status lookups; cache-friendly. \\
$h_D$ oracle & distance + successor arrays & $O(1)$ heuristic query; successor tree enables subtree repair. \\
Obstacle memory & byte array + set & $O(1)$ status; grows monotonically across queries. \\
Repair frontier & set + binary heap & $O((|A|+E_A)\log|A|)$ per discovered obstacle. \\
Risk cache & hash map, versioned & avoids repeated path walks; invalidated on distance changes. \\
\bottomrule
\end{tabular}
\end{table}
